\item \textbf{ACTIVIDAD} Considere los diversos líquidos en la tabla a continuación en sus puntos de ebullición. La tabla muestra el calor latente de vaporización de cada líquido en kJ/mol (note las unidades) y el punto de ebullición en $^\circ\text{C}$. Para cada uno de los líquidos, evalúe el cambio de entropía del líquido por mol cuando se vaporice en el punto de ebullición. ¿Qué advierte sobre los resultados?

\begin{center}
\begin{tabular}{lcc}
\hline
& \textbf{$L_v$ (kJ/mol)} & \textbf{Punto de ebullición ($^\circ$C)} \\
\hline
\textbf{Compuestos polares} & & \\
HF & 25.2 & 19.7 \\
HCl & 16.2 & -84.8 \\
HI & 19.8 & -35.6 \\
H$_2$O & 40.7 & 100 \\
\textbf{Compuestos no polares} & & \\
C$_3$H$_8$ & 19.0 & -42.1 \\
C$_4$H$_{10}$ & 22.4 & -0.50 \\
\textbf{Elementos} & & \\
Hg & 54.7 & 357 \\
Pb & 178 & 1\,749 \\
Cl$_2$ & 20.4 & -34.0 \\
Br$_2$ & 30.0 & 58.8 \\
\hline
\end{tabular}
\end{center}
